\documentclass[11pt,a4paper]{article}
\usepackage[T1]{fontenc}
\usepackage[utf8]{inputenc}
\usepackage[margin=26mm]{geometry}
\usepackage{lmodern,microtype,amsmath,booktabs,tabularx}
\usepackage[hidelinks]{hyperref}
\setlength{\parskip}{0.45em}
\setlength{\parindent}{0pt}
\newcommand{\AllTests}{885}
\newcommand{\LegalCases}{25}
\newcommand{\LegalConditional}{12}
\newcommand{\LegalUnresolved}{13}
\newcommand{\NewRequests}{17}
\newcommand{\RemainingRequests}{25}

\title{When cancellation means a loss\\Prospectus reader repair addendum}
\author{LegalMath development study}
\date{4 October 2026}
\begin{document}
\maketitle

A note can disappear from an issuer's records because the holder has been paid,
or because an unpaid claim has ceased to exist. The same word, \emph{cancelled},
can describe either event. The previous prospectus reader confused these
relations in three clauses from Italian securitisations. This addendum explains
the repair, the opposing repayment controls, and the limits imposed by incomplete
documents and external law. It supplements the preserved thirty-case study; the
original reader, source collection and built study PDFs remain unchanged.

The candidate now retains the unpaid claim, the condition that cancels it and
the exceptions that preserve it. All three original clause checks and both
Mizuho repayment controls pass. The full prospectus suite passes
\AllTests{} tests. These are checks of software behaviour on selected evidence;
independent legal adjudication remains outstanding.

\section*{A claim that survives maturity}

Adriatica Finance's 2007 prospectus supplies the clearest contrast between
maturity and cancellation. Condition 7(b), on PDF page 121, provides that unpaid
amounts continue after the December 2021 maturity when the issuer has insufficient
funds. The claim survives until the later Cancellation Date in December 2026.
At that later date, the cancellation provision retains an exception for the
issuer's gross negligence or wilful misconduct.\footnote{\href{../../prospectus/corner-cases-2026-10-04/sources/jur-more-hypo-843665/source.pdf}{Adriatica Finance, prospectus dated 28 March 2007}, PDF p.~121, Condition 7(b).}

A reader that notices redemption and cancellation but loses this sequence
can erase five years of a surviving claim. A reader that retains the sequence
but omits misconduct can still overstate cancellation. The candidate records
both the cancellation relation and the exact source spans for its prerequisites
and exceptions. Missing facts remain unknown.

For a particular application of this clause, let \(U\) mean that principal is
still unpaid, \(F\) that the relevant insufficiency of funds occurred, and \(D\)
that the Cancellation Date has been reached. Let \(G\) and \(W\) denote issuer
gross negligence and wilful misconduct. Under this bounded reading of the
clause, the cancellation condition is
\[
 C=U\land F\land D\land\neg G\land\neg W.
\]
This expression is a local representation of the reviewed conditions; it does
not decide whether the facts occurred. If the required facts are established
and both exceptions are false, the evaluator returns true. If an exception
holds, it returns false. If a necessary fact remains unknown and no known fact
defeats cancellation, it returns unknown. Tests change each premise separately,
so merely copying the qualifier words into a record is insufficient.

Adriano Lease's 2011 prospectus requires two distinct readings. Condition 8.1.3
cancels specified unpaid amounts at final maturity when funds are insufficient,
except where payment is improperly withheld or refused. Condition 9.2.3 instead
requires a servicer certificate about the absence of reasonable further
recoveries and a representative's notice based on that certificate.\footnote{\href{../../prospectus/corner-cases-2026-10-04/sources/jur-more-hypo-1314069/source.pdf}{Adriano Lease Sec., retained 2011 prospectus}, PDF p.~140, Condition 8.1.3; pp.~144--145, Condition 9.2.3.}
The latter condition crosses a page boundary. The candidate joins the pages
while preserving page locations, and separates neighbouring numbered conditions.
Deleting the certificate, notice or exception from its evidence causes replay
validation to reject that evidence.

These constructions cover the inspected English formulations. They do not
establish that every cancellation clause, definition or cross-reference in an
unseen prospectus has been understood. A conditional positive reading identifies
a documented loss mechanism; it does not establish an actual investor loss.

\section*{A reduction caused by payment}

The opposite mistake appeared in the earlier Mizuho programme study. Its
scheduled-amortisation provisions reduce outstanding principal when the
Amortization Amount is paid. The earlier reader treated that reduction as
principal loss. The two recorded passages are on PDF pages 125--126 and 128
of the 31 August 2026 programme prospectus.\footnote{\href{../../prospectus/difficulty-2026-10-04/pdfs/case-105635287.pdf}{Mizuho programme prospectus}; the original findings and exact quotations remain in the \href{../prospectus-difficulty-2026-10-04/REPORT.md}{preserved difficulty study}.}

Payment explains the change in principal. The candidate therefore checks for
the payment relation before applying a loss label. A separate unpaid
cancellation in the same passage defeats that repayment exclusion. Removing
payment or negating redemption also defeats it. The tests retain both real
Mizuho quotations and these altered controls.

The programme's form fields introduce another limit: an available option has
not necessarily been selected for an issued note. Correcting a clause's
payment meaning does not complete an issue's terms. The fixed-input comparison
therefore preserves the earlier programme scope instead of silently supplying
new final terms.

\begin{center}
\begin{tabularx}{\linewidth}{@{}Xrr@{}}
\toprule
Fixed comparison & Original & Candidate\\
\midrule
Unpaid cancellation wrongly excluded as repayment & 3 & 0\\
Mizuho paid-amortisation loss witnesses in the two controls & 2 & 0\\
Abstentions among seven offering documents & 7 & 5\\
Conditional positive readings among those seven & 0 & 2\\
\bottomrule
\end{tabularx}
\end{center}

The table compares the same retained inputs. The three targeted clause probes
are reported separately in the machine-readable comparison. Neither the lower
abstention count nor the test count measures legal accuracy. Source-complete
issue decisions would require evidence beyond these selected documents.

\section*{Which documents and which dates}

An empty scan cannot establish that a loss clause is absent. The repaired intake
reports two wholly empty extractions and one partial extraction among the seven
offerings, and blocks a definitive contractual answer for those incomplete
inputs. Existing reviewed images remain available, but manual reading is not
presented as automatic OCR. An OCR executable was unavailable on the recorded
launcher path; no installation was made.

BES illustrates why readable text alone is insufficient. Series 23 requires
the 3 November 2010 base prospectus and six supplements; Series 35 and Series 36
require the 17 July 2013 base, with two and four supplements respectively.
The candidate checks the declared issuer, senior or subordinated class,
identifier, document role, edition, base reference and exact required set.
A substituted later base, a missing supplement or the wrong issuer fails that
check. The metadata are source-reviewed declarations; the validator does not
independently infer their meaning from unrestricted prose.

The final-terms scans also distinguish the date of the document from the date
of issue. Page 1 gives 14 July 2011 as the Series 23 final-terms date and
15 July as its issue date; Series 35 uses 20 and 21 January 2014; Series 36 uses
6 and 8 May 2014. The first candidate intake had conflated these dates.
The corrected record retains both dates and the earlier result remains
available for review.\footnote{See \href{SOURCE-REVIEW.md}{source metadata review}, with the precise page and item references.}

All three BES dossiers remain unresolved because the complete incorporated
financial reports and the precedence of the full document set have not been
reviewed. A successful identity check cannot supply an omitted source.
Likewise, a BES Finance guarantee in a shared programme does not give a
BES-issued bond a second independent obligor.

\section*{What external law can change}

A contractual promise to repay principal does not answer whether a later
administrative measure changed the claim, whether the forum recognises that
measure, or whether protected public funds can be attached. The new legal
component evaluates explicitly declared, source-bound premises and keeps these
questions separate. It does not infer a legal rule simply from the identity of
a test case.

The dated BES comparisons distinguish a measure adopted before proceedings
from a retransfer during pending merits proceedings. The Dana Gas scenarios
distinguish an assumption of UAE invalidity from a finding of illegality and
retain the separate English purchase undertaking. The Lloyds comparisons keep
the majority's regulatory-call reasoning separate from dissent and conversion.
The Ukraine scenario prevents permission to try a specified duress defence from
becoming a finding that the debt has been discharged. The retained study supplies
the judgments, page or paragraph anchors and their qualifications.\footnote{\href{../prospectus-corner-cases-2026-10-04/REPORT.md}{Preserved corner-case report} and \href{../prospectus-corner-cases-2026-10-04/test-specifications.json}{thirty source-bound specifications}.}

The additional enforcement scenarios preserve distinctions disclosed in the
Italian prospectuses: a regional receivable is not itself a note guarantee;
pari passu ranking does not establish attachability of protected public funds;
the disclosed 120-day wait is a timing condition; and acceleration of the notes
does not automatically accelerate the regional debt. The day-count evaluator
checks the boundary at 119 and 120 elapsed calendar days under its declared
premises. It does not supply missing service facts or adjudicate legal
calendar-counting rules.

The historical-source search recovered official Italian legislation, including
Article 4 of Law 130 in its version applicable to the requested 2009 snapshot.
Paragraph 3 expressly refers to Article 67 of the bankruptcy law. It does not
establish the separate Article 65 exemption that the 2009 RMBS prospectus
says is absent.\footnote{\href{../../prospectus/corner-repair-2026-10-04/sources/jur-repair-italy-130-art4/source.html}{Normattiva, Law 130, Article 4}, version displayed from 29 May 1999 to 21 February 2014; \href{../../prospectus/corner-cases-2026-10-04/sources/jur-more-hypo-1123807/source.pdf}{Adriano Finance 2009 prospectus}, PDF p.~65.}
Historical Article 67 quoted in explanatory notes is not substituted for the
operative 2009 text. The exact bankruptcy-article requests failed and remain
recorded as source gaps.

All \LegalCases{} declared legal scenarios match their engineering expectations:
\LegalConditional{} produce conditional results and \LegalUnresolved{} remain
unresolved. Changes to claim type, forum, stage, dates, source bytes and
conflicting facts are tested independently. Matching these declared
expectations does not independently validate the legal interpretations from
which they were constructed.

The remaining work is specific: obtain and review complete BES incorporation
and precedence records; recover the Portuguese decision and Annex 2B, the
original HETA measure and liability list, and the missing original offers;
complete the Italian historical-law comparison; then adjudicate the case labels
independently and test unseen formulations. The follow-up used \NewRequests{}
public requests, leaving \RemainingRequests{} of the existing allowance.
The candidate is retained as an isolated, reviewable patch. Human acceptance of
the prose and legal adjudication remain pending.
\end{document}
